%% Owner: Víctor Martín Castilla
\section{High availability and disaster recovery}
\label{sec:hadr}

High availability recovers from minor outages automatically; disaster
recovery restores service after a major one, by hand and from backups
\parencite{schmidt2006ha,llusa2026hadr}.

\subsection{High availability}
\label{sec:hadr-ha}

\begin{itemize}
  \item Three managers keep the Swarm manageable when one node fails
        (\cref{sec:infra-three}).
  \item Stateless services run at least two replicas spread over the
        nodes. Swarm restarts failed tasks elsewhere and the routing mesh
        balances requests between replicas.
  \item Health checks let Swarm replace containers that are running but
        no longer working \parencite{llusa2026production}.
  \item Rolling updates replace one replica at a time and roll back
        automatically when the new version fails.
\end{itemize}

\todo{List the single points of failure of the cluster the team uses,
such as the public entry, the NFS server, the database node and the
Azure region, and measure what happens when each node is stopped.}

\subsection{Disaster recovery}
\label{sec:hadr-dr}

\todo{choose the service category with the team and justify it. As a
starting point, Schmidt's ``business-important'' category allows a
recovery time objective (RTO) of three days and a recovery point
objective (RPO) of one day \parencite{schmidt2006ha}, which daily dumps
satisfy. Document the recovery procedure from zero and its measured
duration.}
